\documentclass[aspectratio=169,11pt]{beamer}

\usetheme{Madrid}
\usecolortheme{default}
\usefonttheme{professionalfonts}
\setbeamertemplate{navigation symbols}{}
\setbeamertemplate{footline}[frame number]

\usepackage{amsmath, amssymb, booktabs, array}

\title{Synthesis, Frontier Questions, And Student Research Designs}
\subtitle{Gender Study --- Week 10}
\author{Teaching deck draft}
\date{}

\begin{document}

\begin{frame}
  \titlepage
\end{frame}

\begin{frame}{Central question}
\begin{itemize}
    \item What does a credible frontier research design on gender and labor markets look like?
    \item How do we move from a gendered labor fact to a mechanism, counterfactual, method, and contribution?
    \item Why should a setting-specific project be broadly interesting to labor economists?
\end{itemize}
\vspace{0.25cm}
\[
\text{labor object} \rightarrow \text{mechanism} \rightarrow
\text{counterfactual} \rightarrow \text{identified margin} \rightarrow
\text{contribution}
\]
\end{frame}

\begin{frame}{Course recap as a research map}
\scriptsize
\begin{tabular}{p{0.17\linewidth} p{0.35\linewidth} p{0.36\linewidth}}
\toprule
Block & Labor objects & Research-design question \\
\midrule
Week 1 & gaps, measurement, welfare & what outcome is being measured? \\
Weeks 2--3 & care, fertility, skills, sorting & when and why do trajectories diverge? \\
Weeks 4--5 & firms, authority, norms, bargaining & what is worker sorting versus treatment? \\
Weeks 6--7 & policy, safety, mobility & who bears incidence and hidden harms? \\
Weeks 8--9 & comparative settings, new data & what travels, and what can data see? \\
Optional & LGBTQ+ labor institutions & what identity object is observed? \\
Week 10 & student research designs & what is the contribution? \\
\bottomrule
\end{tabular}
\end{frame}

\begin{frame}{Where the literature is now}
\begin{itemize}
    \item Beyond a single wage gap: households, firms, policy, safety, institutions, identity, and welfare.
    \item Lifecycle evidence made care, children, and dynamic careers central.
    \item Firm-side work opened hiring, promotions, pay-setting, networks, authority, and job design.
    \item Policy papers increasingly ask incidence: workers, firms, households, applicants, coworkers.
    \item New data reveal applications, postings, managers, time use, platform choices, climate, and identity.
\end{itemize}
\end{frame}

\begin{frame}{What has been answered}
\begin{itemize}
    \item Children and care are central to persistent earnings, hours, and career gaps.
    \item Firms matter: sorting across firms and treatment within firms both shape gaps.
    \item Policy labels are not mechanisms; financing, compliance, take-up, and firm response matter.
    \item Comparative evidence is useful when it isolates mechanisms, not when it tours countries.
    \item Measurement is part of the estimand: legal category, self-identity, perceived identity, and disclosure differ.
\end{itemize}
\end{frame}

\begin{frame}{Unresolved core debates and gaps}
\scriptsize
\begin{tabular}{p{0.29\linewidth} p{0.55\linewidth}}
\toprule
Frontier gap & Contribution margin \\
\midrule
Firm organization & authority, evaluation, schedules, managers, low-promotability tasks \\
Hidden harms & safety, harassment, mobility risk, stigma, dignity, mental health \\
Identity measurement & LGBTQ+, trans, nonbinary, disclosure, perceived treatment, privacy \\
Policy incidence & firm response, applicants, households, nonparticipants, coworkers \\
Portability & what travels: mechanism, sign, magnitude, policy, or method? \\
Technology & AI exposure, monitoring, scheduling, task change, remote work \\
\bottomrule
\end{tabular}
\end{frame}

\begin{frame}{Research-design template}
\scriptsize
\begin{tabular}{p{0.28\linewidth} p{0.55\linewidth}}
\toprule
Step & Question \\
\midrule
Labor object & what outcome or allocation problem is central? \\
Mechanism & why should gender shape the object? \\
Counterfactual & compared with whom, what, or when? \\
Labor margin & search, hiring, hours, promotion, firm sorting, safety, benefits? \\
Data source & what data observe the moving margin? \\
Identification & what variation distinguishes mechanism from selection? \\
Welfare object & whose income, time, risk, autonomy, authority, or dignity changes? \\
Contribution & what should readers update about labor economics? \\
\bottomrule
\end{tabular}
\end{frame}

\begin{frame}{From topic to design}
\begin{columns}[T,onlytextwidth]
\begin{column}{0.47\linewidth}
\textbf{Weak start}
\begin{itemize}
    \item ``Gender and commuting''
    \item ``Women in firms''
    \item ``LGBTQ+ wage gaps''
    \item ``Childcare policy''
\end{itemize}
\end{column}
\begin{column}{0.50\linewidth}
\textbf{Research-design start}
\begin{itemize}
    \item safety risk and search radius
    \item manager assignment and promotion ratings
    \item identity disclosure and firm sorting
    \item childcare access and retention at high-premium firms
\end{itemize}
\end{column}
\end{columns}
\vspace{0.25cm}
The design begins when the moving labor margin is visible.
\end{frame}

\begin{frame}{Identification threats to diagnose early}
\begin{itemize}
    \item Selection into observed work, firms, applications, surveys, or reports.
    \item Controls that are mechanisms: occupation, hours, tenure, firm, manager, promotion.
    \item Anticipation before childbirth, policy rollout, disclosure, reporting, or legal change.
    \item Reporting versus prevalence in complaints, harassment, safety, and climate data.
    \item Firm response: pay, hiring, job design, outsourcing, sorting, compliance.
    \item Measurement mismatch between the available category and the treatment concept.
\end{itemize}
\end{frame}

\begin{frame}{Portability and broad-interest framing}
\scriptsize
\begin{tabular}{p{0.24\linewidth} p{0.58\linewidth}}
\toprule
Question & Strong answer \\
\midrule
Why this setting? & it isolates a labor mechanism unusually clearly \\
What travels? & the mechanism or margin, even if levels differ \\
What does not travel? & local legal, cultural, sectoral, or enforcement details \\
Why should nonexperts care? & the result informs a broader labor-market question \\
What is the contribution? & work, wages, firms, households, mobility, institutions, or welfare \\
\bottomrule
\end{tabular}
\vspace{0.25cm}
The setting is leverage; the contribution is the labor mechanism.
\end{frame}

\begin{frame}{Interdisciplinary directions}
\scriptsize
\begin{tabular}{p{0.28\linewidth} p{0.58\linewidth}}
\toprule
Direction & Labor question \\
\midrule
Urban and spatial & commuting, housing, safety, childcare geography, job access \\
Political economy & representation, enforcement, public employment, voice, law \\
Health and reproduction & fertility timing, IVF, abortion access, menopause, well-being \\
Technology and AI & task exposure, monitoring, scheduling, remote work, platforms \\
Crime and safety & violence, harassment, mobility, sorting, welfare \\
Organizations & leadership, mentoring, culture, evaluation, authority \\
\bottomrule
\end{tabular}
\end{frame}

\begin{frame}{Project-generation guidance}
\begin{itemize}
    \item Start with one labor object.
    \item Choose one mechanism that could plausibly move it.
    \item Name the comparison group before choosing the estimator.
    \item Match the data to the moving margin.
    \item Write the main threat in one sentence.
    \item State what travels beyond the setting.
    \item End with the welfare or distributional object.
\end{itemize}
\end{frame}

\begin{frame}{Lab: Reproduce, Diagnose, Design}
\begin{columns}[T,onlytextwidth]
\begin{column}{0.32\linewidth}
\textbf{Reproduce}
\begin{itemize}
    \item build idea bank
    \item classify objects
    \item export design table
\end{itemize}
\end{column}
\begin{column}{0.32\linewidth}
\textbf{Diagnose}
\begin{itemize}
    \item score fit
    \item flag threats
    \item compare methods
\end{itemize}
\end{column}
\begin{column}{0.32\linewidth}
\textbf{Design}
\begin{itemize}
    \item pick one idea
    \item memo scaffold
    \item contribution sentence
\end{itemize}
\end{column}
\end{columns}
\end{frame}

\begin{frame}{Bridge to students' own agendas}
\begin{itemize}
    \item A good gender-and-labor project is precise about mechanism.
    \item It is honest about what the design identifies and what remains hidden.
    \item It treats local knowledge as leverage, not as the whole contribution.
    \item It makes welfare broader than wages when safety, time, dignity, autonomy, or authority matter.
    \item The goal is not to add gender to a labor topic; it is to learn labor economics through gendered mechanisms.
\end{itemize}
\end{frame}

\begin{frame}{Final takeaway}
\Large
\[
\text{fact} \neq \text{mechanism} \neq \text{design} \neq \text{contribution}
\]
\vspace{0.4cm}
\normalsize
Week 10 is the handoff: from course map to research agenda.
\end{frame}

\end{document}
